% sec-04: authority boundaries.  Figure 2 is a plantuml-type use case diagram.
\section{Who Is Authorized to Do What}

The program has three authorities that must not blur: the operating surgeon
holds clinical authority, the program manager holds programmatic authority, and
the advisory model holds none. Writing that down precisely is not a formality,
because the single most common objection to an LLM in an operating room is that
the boundary is asserted rather than specified.

\begin{appfloat}[!tb]
\begin{appfig}[plantuml-type / use case]
% Three actors on the left, one on the right, use cases in a bounded system box.
% Actors are 1.7 apart vertically; ellipses are 1.35 apart, so no association
% line crosses an ellipse.
\node[umlpkg,minimum width=6.6cm,minimum height=5.5cm] (sys) at (4.6,-1.35) {};
\node[umlpkgtab,anchor=north west] at ([xshift=1mm,yshift=-1mm]sys.north west) {Trial system boundary};
\node[umlusekey] (u1) at (4.6,0.05) {Approve every robotic motion};
\node[umlusecase] (u2) at (4.6,-1.3) {Order conversion to open operation};
\node[umlusecase] (u3) at (4.6,-2.65) {Render advisory text to a display};
\umlactor{a1}{0}{0.15}{Operating surgeon}
\umlactor{a2}{0}{-2.0}{ARPA-H program manager}
\umlactor{a3}{9.4}{-2.55}{On-premises advisory model}
\draw[umlassoc] (a1) -- (u1.west);
\draw[umlassoc] (a1) -- (u2.west);
\draw[umlassoc] (a2) -- node[umlguard,pos=0.55] {gate decisions only} (u2.west);
\draw[umlassoc] (a3) -- (u3.east);
\draw[umldash] (u3.north) to[out=100,in=-80,looseness=0.8] node[umlguard,pos=0.5] {advisory input} (u1.south);
\node[umlnote,anchor=north west,text width=42mm] (n1) at (7.8,0.55)
  {The model has no association with the two authorized actions. No generated
  token reaches a motion controller.};
\draw[umldash] (n1.west) -- (u1.east);
\pxmark{7.35}{-2.55}
\node[font=\tiny\sffamily,text=protogray,anchor=north,text width=34mm,align=center]
  at (7.35,-2.95) {no path to the arms};
\end{appfig}
\vspace{-0.7cm}
\figcaption{Authority as an association diagram. Only the surgeon is associated
with the two authorized actions; the advisory model is associated with one
action, and the struck link is the one it does not have.}
\end{appfloat}

The operation itself is a staged eight-arm robotic pancreaticoduodenectomy with
perioperative daraxonrasib \cite{onpremwhippl,phase1ind}. Arm-level stop is
specified at 3 milliseconds and system-wide stop at 500 milliseconds, both
measured on the bench before the first patient and re-measured at every cohort
boundary. Conversion ordered for patient safety counts as clinically
appropriate; conversion caused by device, advisory, or integration failure
counts against the feasibility endpoint and is reported separately.

\begin{apptable}
\begin{tabularx}{\textwidth}{>{\raggedright\arraybackslash}p{3.5cm} >{\raggedright\arraybackslash}p{3.0cm} >{\raggedright\arraybackslash}X}
\headrow \thead{Surgical measure} & \thead{Recorded as} & \thead{Gate that reads it}\\
\midrule
Protocol tasks completed without unsafe conversion & Count and proportion, with attribution & Gate 2 and Gate 3 \\
Operative time and estimated blood loss & Minutes and millilitres per case & Gate 3, trend only \\
Pancreatic fistula, ISGPS grade & Grade B or C at 90 days \cite{bassi2017isgps} & Gate 3 \\
R0 margin status & Proportion of resections & Gate 3, supportive \\
\end{tabularx}
\end{apptable}
\tabcap{Four surgical measures, how each is recorded, and which gate reads it.
Attribution is recorded at the time of the event rather than reconstructed at
the gate.}

\vspace{0.1em}
\noindent\footnotesize\textbf{Positioning note.} \textit{First in human refers
to the integrated surgical and advisory workflow, not to daraxonrasib, which is
investigational and already in Phase 3 evaluation. The robotic configuration is
specified by manufacturer, model, cleared configuration, arm count, and software
version at the site agreement, and no drug supply or letter of authorization is
in place with the agent's developer.}\normalsize\par
